\documentclass[10pt,a4paper]{amsart}
\usepackage[margin=19mm]{geometry}
\usepackage{amsmath,amssymb,mathtools}
\usepackage[scr=rsfs,cal=euler]{mathalfa}
\usepackage{xfrac}
\usepackage[unicode,hidelinks,bookmarks=false]{hyperref}
\newcommand{\ie}{i.e.\ }
\newcommand{\cf}{cf.\ }
\newcommand{\resp}{respectively\ }
\newcommand{\Subsection}{\subsection}
\providecommand{\nobreakdash}{\nobreak\hbox{-}}
\setlength{\emergencystretch}{2em}
\multlinegap0pt
\usepackage{algorithm}\usepackage{algpseudocode}
\usepackage{tikz}
\usepackage{amssymb}
\newtheorem{proposition}{Proposition}
\title{Secrecy Outage Analysis over Correlated Composite Generalized-Gamma Fading Channels}
\author{Muhammad Adil, Pranavkumar Pathak,, Salman A. Alqahtani}
\date{Source: arXiv:2608.29414v1 (CC BY 4.0)}
\begin{document}
\maketitle

\noindent\emph{Source and licence.} Secrecy Outage Analysis over Correlated Composite Generalized-Gamma Fading Channels. Version arXiv:2608.29414v1 (CC BY 4.0), distributed by arXiv under the Creative Commons Attribution 4.0 International licence, http://creativecommons.org/licenses/by/4.0/, as recorded in the arXiv licence metadata for this exact version and shown on the abstract page https://arxiv.org/abs/2608.29414v1. Full licence notice: \texttt{licenses/arxiv-2608.29414-CC-BY-4.0.txt}.\par\smallskip

\noindent\textit{Editorial scope.} Counted unit: the article's proposition prop:4prime (source0.tex lines 410--465). The article's complete original proof of it is delivered with it (source0.tex lines 467--508, one \texttt{proof} environment of 42 lines). No mathematics is written, repaired or re-derived: the statement and the proof are the article's own lines, in the article's own notation, and no retained line is substituted or re-flowed.  Retained uncounted as the source-local context the proof needs: lines 208--215; lines 281--291.  Nothing was omitted from any retained span, so the package carries no omission record.  No retained line contains a citation, so the wrapper carries no bibliography and no bibliography entry.  No retained line contains a figure environment, an image-inclusion command or drawing code, so no image is delivered and the package declares no asset dependency.  Editorial formatting only, in addition: an AMS \texttt{amsart} preamble that restates the article's own declarations and macros used by the retained lines, namely the environments \texttt{proposition} on separate counters, together with the standard packages those lines need.  The retained environments print in this document as Proposition 1 in order of appearance, whereas the article prints the same items with the numbering of its own full document.  The complete native source of the article as distributed in the arXiv e-print for this exact version is bundled unchanged as \texttt{sources/208851/source0.tex}.
\begin{equation}
\mathcal M_Z(\zeta)
=D^{\zeta-1}
\frac{\Gamma\!\left(k+\frac{\zeta-1}{\beta}\right)
\Gamma\!\left(m+\frac{\zeta-1}{\kappa}\right)}
{\Gamma(k)\Gamma(m)},
\label{eq:MZ}
\end{equation}

\begin{equation}
w_{ij}
=
\frac{(1-\rho)^{k_{\mathrm E}}
(k_{\mathrm B})_i
(k_{\mathrm E}-k_{\mathrm B})_j
\rho^{i+j}}
{i!\,j!},
\qquad i,j\in\mathbb N_0.
\label{eq:corrected_weights}
\end{equation}

\begin{proposition}\label{prop:4prime}
For the complex Mellin variable \(q\) satisfying \(\Re(q)>0\), the Mellin transform of the legitimate-link survival function is
\begin{equation}
\mathcal M_{G_{\mathrm B,ij}}(q)
=
\frac{D_{\mathrm B}^q}{q}\,
\frac{\Gamma\!\left(m_{\mathrm B}+\frac{q}{\kappa_{\mathrm B}}\right)
\Gamma\!\left(s_{\mathrm B}+\frac{q}{\beta_{\mathrm B}}\right)}
{\Gamma(m_{\mathrm B})\Gamma(s_{\mathrm B})}.
\label{eq:MG1prime}
\end{equation}
For \(\Re(q)>1-\min\{\kappa_{\mathrm E}m_{\mathrm E},
\beta_{\mathrm E}s_{\mathrm E}\}\), the Mellin transform of the eavesdropping-link PDF is
\begin{equation}
\mathcal M_{f_{\mathrm E,ij}}(q)
=
D_{\mathrm E}^{q-1}\,
\frac{\Gamma\!\left(m_{\mathrm E}+\frac{q-1}{\kappa_{\mathrm E}}\right)
\Gamma\!\left(s_{\mathrm E}+\frac{q-1}{\beta_{\mathrm E}}\right)}
{\Gamma(m_{\mathrm E})\Gamma(s_{\mathrm E})}.
\label{eq:Mkappa2prime}
\end{equation}
Using Mellin--Parseval's identity over a vertical contour \(\mathcal L\) satisfying \(0<\Re(q)<\min\{\kappa_{\mathrm E}m_{\mathrm E},\beta_{\mathrm E}s_{\mathrm E}\}\),
\[
Q_{ij}
=
\frac{1}{2\pi\mathrm i}
\int_{\mathcal L}
\mathcal M_{G_{\mathrm B,ij}}(q)\,
\mathcal M_{f_{\mathrm E,ij}}(1-q)\,dq .
\]
Define \(\mathbf a_{ij}\triangleq\left\{\left(1-m_{\mathrm E},\frac{1}{\kappa_{\mathrm E}}\right),\left(1-s_{\mathrm E},\frac{1}{\beta_{\mathrm E}}\right),(1,1)\right\}\) and \(\mathbf b_{ij}\triangleq\left\{\left(m_{\mathrm B},\frac{1}{\kappa_{\mathrm B}}\right),\left(s_{\mathrm B},\frac{1}{\beta_{\mathrm B}}\right),(0,1)\right\}\).
Then,
\begin{equation}
Q_{ij}
=
\frac{1}
{\Gamma(s_{\mathrm B})\Gamma(m_{\mathrm B})
\Gamma(s_{\mathrm E})\Gamma(m_{\mathrm E})}
H_{3,3}^{3,2}\!\left[
\frac{D_{\mathrm E}}{D_{\mathrm B}}
\;\middle|\;
\mathbf a_{ij}\,;\mathbf b_{ij}
\right].
\label{eq:Qijprime}
\end{equation}
Consequently,
\begin{equation}
\mathrm{PNZSC}
=
\sum_{i=0}^{\infty}\sum_{j=0}^{\infty}w_{ij}Q_{ij},
\qquad
P_o(0)=1-\mathrm{PNZSC}.
\label{eq:PNZSCprime}
\end{equation}
\end{proposition}

\begin{proof}
The Mellin transform in \eqref{eq:MG1prime} follows from the general tail identity
\[
\int_0^\infty x^{q-1}\Pr[Z>x]\,dx
=
\frac{\mathbb E[Z^q]}{q},
\qquad \Re(q)>0,
\]
together with the \(q\)-th moment obtained from \eqref{eq:MZ}. Equation~\eqref{eq:Mkappa2prime} follows directly from \eqref{eq:MZ} after replacing \(k\), \(m\), \(\beta\), \(\kappa\), and \(D\) by \(s_{\mathrm E}\), \(m_{\mathrm E}\), \(\beta_{\mathrm E}\), \(\kappa_{\mathrm E}\), and \(D_{\mathrm E}\), respectively.

Substituting \eqref{eq:MG1prime} and \eqref{eq:Mkappa2prime} into the Parseval integral gives
\begin{equation}
\begin{aligned}
&Q_{ij}\Gamma(s_{\mathrm B})\Gamma(m_{\mathrm B})
\Gamma(s_{\mathrm E})\Gamma(m_{\mathrm E})\\
&=
\frac{1}{2\pi\mathrm i}
\int_{\mathcal L}
\frac{
\Gamma\!\left(m_{\mathrm B}+\frac{q}{\kappa_{\mathrm B}}\right)
\Gamma\!\left(s_{\mathrm B}+\frac{q}{\beta_{\mathrm B}}\right)}
{q}\\
&\quad\times
\Gamma\!\left(m_{\mathrm E}-\frac{q}{\kappa_{\mathrm E}}\right)
\Gamma\!\left(s_{\mathrm E}-\frac{q}{\beta_{\mathrm E}}\right)
\left(\frac{D_{\mathrm B}}{D_{\mathrm E}}\right)^q dq .
\end{aligned}
\label{eq:Qij_MB_integral}
\end{equation}
Using \(\frac{1}{q}=\frac{\Gamma(q)}{\Gamma(1+q)}\) and \(\left(\frac{D_{\mathrm B}}{D_{\mathrm E}}\right)^q=\left(\frac{D_{\mathrm E}}{D_{\mathrm B}}\right)^{-q}\), the factors in \eqref{eq:Qij_MB_integral} are identified with the standard Mellin--Barnes definition of the Fox \(H\)-function as follows. By comparison with the standard Mellin--Barnes representation, the Fox-\(H\) orders are \(p_{\mathrm H}=q_{\mathrm H}=m_{\mathrm H}=3\) and \(n_{\mathrm H}=2\), with
\[
\begin{gathered}
(a_1,A_1)=(1-m_{\mathrm E},\kappa_{\mathrm E}^{-1}),\quad
(a_2,A_2)=(1-s_{\mathrm E},\beta_{\mathrm E}^{-1}),\\
(a_3,A_3)=(1,1),\quad
(b_1,B_1)=(m_{\mathrm B},\kappa_{\mathrm B}^{-1}),\\
(b_2,B_2)=(s_{\mathrm B},\beta_{\mathrm B}^{-1}),\quad
(b_3,B_3)=(0,1).
\end{gathered}
\]
Collecting two reflected upper-numerator terms, three direct lower-numerator terms, one direct upper-denominator term, and no reflected lower-denominator terms yields the Fox \(H\)-function in \eqref{eq:Qijprime}. Summing over \((i,j)\) with the normalized weights in \eqref{eq:corrected_weights} gives \eqref{eq:PNZSCprime}.
\end{proof}

\end{document}
